% Appendix: services and compositions of the evaluation (referenced from Sec. V, Table~\ref{tab:applications}).
% Sources: service images (spec.yaml, code) and statics/results_milos_v2/{benchmark/w1,b1,b1-unthrottled}; the compositions are
% the chains of DEFAULT_CHAINS (b1 + b1u-max4) in notebooks/tsc/pegasus/evaluation_setup.py.
%
% Usage in main.tex, after \appendices:  % Appendix: services and compositions of the evaluation (referenced from Sec. V, Table~\ref{tab:applications}).
% Sources: service images (spec.yaml, code) and statics/results_milos_v2/{benchmark/w1,b1,b1-unthrottled}; the compositions are
% the chains of DEFAULT_CHAINS (b1 + b1u-max4) in notebooks/tsc/pegasus/evaluation_setup.py.
%
% Usage in main.tex, after \appendices:  % Appendix: services and compositions of the evaluation (referenced from Sec. V, Table~\ref{tab:applications}).
% Sources: service images (spec.yaml, code) and statics/results_milos_v2/{benchmark/w1,b1,b1-unthrottled}; the compositions are
% the chains of DEFAULT_CHAINS (b1 + b1u-max4) in notebooks/tsc/pegasus/evaluation_setup.py.
%
% Usage in main.tex, after \appendices:  % Appendix: services and compositions of the evaluation (referenced from Sec. V, Table~\ref{tab:applications}).
% Sources: service images (spec.yaml, code) and statics/results_milos_v2/{benchmark/w1,b1,b1-unthrottled}; the compositions are
% the chains of DEFAULT_CHAINS (b1 + b1u-max4) in notebooks/tsc/pegasus/evaluation_setup.py.
%
% Usage in main.tex, after \appendices:  \input{appendix_services}
% \twocolumn[...] starts a new page and sets the heading and the service table across both columns at its top;
% the text below continues in two columns. The table is no float (a table* would be pushed to the page after),
% so \@captype gives it a caption and number of a regular table.
% Service names and the order of the applications (Audio, Lidar, Vision) follow Table~\ref{tab:applications}.
% The text only adds what Sec. V does not state; it must fit on one page together with the table.
% The footnotes use \url (packages url and hyperref in main.tex).

\makeatletter
\twocolumn[{%
\section{Detailed Description of Services and Compositions}
\label{apx:services}
\begin{center}
    \def\@captype{table}
    \caption{Services per application. \textit{Fixed knobs} stay at their default in all runs; \textit{Data Q.}\ is the data quality range covered by the benchmarks, \textit{Proc.}\ the range of the median processing time of an item over the benchmarked configurations with one worker, and \textit{Runs} the number of valid benchmark runs with 1--4 workers.}
    \label{tab:apx-services}
    \scriptsize
    \setlength{\tabcolsep}{4pt}
    \begin{tabular}{@{}l p{0.33\textwidth} l l l l r r@{}}
        \toprule
        {Service} & {Function} & {Output} & {Built on} & {Fixed knobs} & {Data Q.} & {Proc.\ [ms]} & {Runs} \\
        \midrule
        \multicolumn{8}{@{}l}{\textit{Audio} (data quality: clip length, seconds)} \\
        \textbf{resample} & resamples a clip, removes its offset, normalizes loudness & audio & NumPy, SciPy & target rate & 2--12 & 0.6--3.9 & 200 \\
        \textbf{denoise} & suppresses background noise in the spectrum & audio & NumPy & FFT size & 2--12 & 6.4--35 & 194 \\
        \textbf{MFCC} & computes a mel spectrogram and its MFCC features & features & NumPy & mel bands, frame & 2--12 & 0.9--4.7 & 200 \\
        \textbf{voicedetect} & returns the start and end times of speech segments & JSON & NumPy & frame length & 3--12 & 4.3--16 & 200 \\
        \textbf{pitch} & estimates the pitch of the voice per frame & pitch & NumPy & hop length & 2--12 & 7.6--31 & 187 \\
        \textbf{transcribe} & transcribes speech to text (small model) & text & Vosk & model size & 2--12 & 337--983 & 164 \\
        \textbf{encode} & encodes a clip as a compressed audio file & bytes & NumPy, soundfile & format & 2--12 & 1.6--9.1 & 200 \\
        \midrule
        \multicolumn{8}{@{}l}{\textit{Lidar} (data quality: max.\ distance, meters)} \\
        \textbf{voxel} & downsamples a scan to one point per voxel & points & NumPy & voxel size & 8--40 & 1.9--9.0 & 197 \\
        \textbf{ground} & fits the ground plane by random trials, removes it & points & NumPy & attempts & 8--40 & 17--47 & 194 \\
        \textbf{cluster} & clusters the points, returns one box per object & JSON & NumPy, SciPy & voxel size, distance & 6--40 & 0.2--16 & 199 \\
        \textbf{normals} & estimates the surface normal of every point & normals & NumPy, SciPy & neighbors & 6--40 & 1.8--46 & 198 \\
        \textbf{register} & aligns a scan with the previous one (rotation, translation) & JSON & NumPy, SciPy & iterations & 6--38 & 1.0--45 & 189 \\
        \textbf{birdview} & renders a bird's-eye-view image of the scan & image & NumPy & -- & 6--40 & 1.1--3.2 & 197 \\
        \textbf{encode} & quantizes and packs the point coordinates & bytes & NumPy & precision & 6--38 & 3.2--18 & 200 \\
        \midrule
        \multicolumn{8}{@{}l}{\textit{Vision} (data quality: frame height, pixels; for object and face also the model size)} \\
        \textbf{object} & detects objects (YOLOv10): boxes, classes, scores & JSON & ONNX Runtime & -- & 160--640; 1--2 & 13--320 & 182 \\
        \textbf{face} & detects faces (YuNet): boxes and scores & JSON & OpenCV & -- & 256--640; 1--2 & 7.6--195 & 171 \\
        \textbf{edge} & finds edges (bilateral filter, Canny) and their contours & JSON & OpenCV & smoothing & 160--640 & 0.9--6.2 & 200 \\
        \textbf{motion} & returns the regions moved since the previous frame & JSON & OpenCV & method & 160--640 & 0.8--6.0 & 200 \\
        \textbf{QR} & finds and decodes QR codes, draws them onto the frame & image & OpenCV, pyzbar & -- & 192--608 & 1.5--9.1 & 200 \\
        \textbf{transform} & resizes a frame, corrects its colors and geometry & image & OpenCV & correction & 192--640 & 7.1--41 & 200 \\
        \textbf{compress} & encodes a frame as a compressed image & bytes & OpenCV, Pillow & quality, format & 160--608 & 0.2--1.0 & 181 \\
        \textbf{denoise} & removes noise, keeps edges (bilateral filter) & image & OpenCV & strength & 160--640 & 1.4--22 & 198 \\
        \bottomrule
    \end{tabular}
\end{center}
\vspace{6pt}
}]
\makeatother

\noindent
\textbf{Services.}
%
Items are LibriSpeech \textit{dev-clean} clips of at least 12\,s\footnote{\url{https://www.openslr.org/12}}, scans of the KITTI raw recording \texttt{2011\_09\_26\_drive\_0001\_sync}\footnote{\url{https://www.cvlibs.net/datasets/kitti/raw_data.php}}, and frames of sample videos. The data quality knob sets the length of the processed clip, the radius around the car within which points are kept, and the height to which a frame is resized. Three services run pretrained models: transcribe (Vosk\footnote{\url{https://alphacephei.com/vosk/models}}), object (YOLOv10\footnote{\url{https://github.com/THU-MIG/yolov10}}), and face (YuNet\footnote{\url{https://github.com/opencv/opencv_zoo/tree/main/models/face_detection_yunet}}). For object and face, the model size is a second data quality knob: 1 selects a small model, 2 a larger one, which takes several times as long per frame. All other knobs stay at their default.

\vspace{3pt}
\noindent
\textbf{Benchmarks.}
%
Every service runs alone with $c = 1, \dots, 4$ workers in the same 10 configurations, each at the utilizations 0.1, 0.4, 0.6, 0.8, and 1.0 per worker. A run has the CPU need of its configuration, measured in an 8\,s probe, plus 15\% as cores, and measures 30\,s or 400 items (at most 180\,s) after a 10\,s warm-up. Invalid runs, in which the service fell behind its arrivals or needed more cores than its probe had, are dropped (\textit{Runs}). The GPs of a service take as inputs its data quality (and model size), its workers $c$, and the logarithm of its arrival rate in items/s; the cores are no input, since no used run is throttled. Where a benchmark covers less data quality than the compositions (e.g., face from 256\,px), the models extrapolate.

\vspace{3pt}
\noindent
\textbf{Compositions.}
%
Each composition is measured in 200 runs with the measurement window of the benchmarks: 150 runs with one worker and 1, 1.5, or 2 cores per service, and 50 runs with $c \in \{1, \dots, 4\}$ workers and $c + 0.5$ cores per service. The data quality is drawn in steps of 1\,s (Audio), 2\,m (Lidar), and 32\,px (Vision). Only runs in which no service was throttled are used: 125--156 per composition (75--106 with one worker, all 50 with several).

\vspace{3pt}
\noindent
\textbf{Example.}
%
% Source: statics/results_milos_v2/b1/b1-vision-5-single/*/runs/run-0000 (trace_steps.csv, <node>/records.csv, run.json);
% the item motion-376, the median end-to-end latency among the 415 items of the measurement phase.
Table~\ref{tab:apx-trace} follows one item through the six services of V5 in a run with 576\,px, model size 1, and one worker per service; its end-to-end latency is the median of the 415 items of that run. Four services return JSON, QR the annotated frame as a raw image, and compress a JPEG (quality 75). The 157\,ms of the item are the sum of the times in the table; processing at the object detector alone takes two thirds of it.

\makeatletter
\begin{center}
    \def\@captype{table}
    \caption{One item of V5 per step (576\,px, model size 1): output and its size, time waited in the queue, handed over from the previous service, and processed, in ms.}
    \label{tab:apx-trace}
    \footnotesize
    \setlength{\tabcolsep}{4pt}
    \begin{tabular}{@{}l l r r r r@{}}
        \toprule
        {Service} & {Output} & {Size} & {Wait} & {Handover} & {Proc.} \\
        \midrule
        \textbf{motion} & JSON (moved regions) & -- & 0.3 & -- & 5.6 \\
        \textbf{QR} & raw image (annotated) & 1.77\,MB & 0.4 & 1.9 & 8.1 \\
        \textbf{object} & JSON (detections) & -- & 3.8 & 1.3 & 102.7 \\
        \textbf{face} & JSON (faces) & -- & 0.4 & 1.5 & 21.2 \\
        \textbf{edge} & JSON (contours) & -- & 0.4 & 1.5 & 5.6 \\
        \textbf{compress} & JPEG & 38.8\,kB & 0.4 & 1.3 & 1.0 \\
        \midrule
        \multicolumn{5}{@{}l}{end-to-end latency} & 157.3 \\
        \bottomrule
    \end{tabular}
\end{center}
\makeatother

% \twocolumn[...] starts a new page and sets the heading and the service table across both columns at its top;
% the text below continues in two columns. The table is no float (a table* would be pushed to the page after),
% so \@captype gives it a caption and number of a regular table.
% Service names and the order of the applications (Audio, Lidar, Vision) follow Table~\ref{tab:applications}.
% The text only adds what Sec. V does not state; it must fit on one page together with the table.
% The footnotes use \url (packages url and hyperref in main.tex).

\makeatletter
\twocolumn[{%
\section{Detailed Description of Services and Compositions}
\label{apx:services}
\begin{center}
    \def\@captype{table}
    \caption{Services per application. \textit{Fixed knobs} stay at their default in all runs; \textit{Data Q.}\ is the data quality range covered by the benchmarks, \textit{Proc.}\ the range of the median processing time of an item over the benchmarked configurations with one worker, and \textit{Runs} the number of valid benchmark runs with 1--4 workers.}
    \label{tab:apx-services}
    \scriptsize
    \setlength{\tabcolsep}{4pt}
    \begin{tabular}{@{}l p{0.33\textwidth} l l l l r r@{}}
        \toprule
        {Service} & {Function} & {Output} & {Built on} & {Fixed knobs} & {Data Q.} & {Proc.\ [ms]} & {Runs} \\
        \midrule
        \multicolumn{8}{@{}l}{\textit{Audio} (data quality: clip length, seconds)} \\
        \textbf{resample} & resamples a clip, removes its offset, normalizes loudness & audio & NumPy, SciPy & target rate & 2--12 & 0.6--3.9 & 200 \\
        \textbf{denoise} & suppresses background noise in the spectrum & audio & NumPy & FFT size & 2--12 & 6.4--35 & 194 \\
        \textbf{MFCC} & computes a mel spectrogram and its MFCC features & features & NumPy & mel bands, frame & 2--12 & 0.9--4.7 & 200 \\
        \textbf{voicedetect} & returns the start and end times of speech segments & JSON & NumPy & frame length & 3--12 & 4.3--16 & 200 \\
        \textbf{pitch} & estimates the pitch of the voice per frame & pitch & NumPy & hop length & 2--12 & 7.6--31 & 187 \\
        \textbf{transcribe} & transcribes speech to text (small model) & text & Vosk & model size & 2--12 & 337--983 & 164 \\
        \textbf{encode} & encodes a clip as a compressed audio file & bytes & NumPy, soundfile & format & 2--12 & 1.6--9.1 & 200 \\
        \midrule
        \multicolumn{8}{@{}l}{\textit{Lidar} (data quality: max.\ distance, meters)} \\
        \textbf{voxel} & downsamples a scan to one point per voxel & points & NumPy & voxel size & 8--40 & 1.9--9.0 & 197 \\
        \textbf{ground} & fits the ground plane by random trials, removes it & points & NumPy & attempts & 8--40 & 17--47 & 194 \\
        \textbf{cluster} & clusters the points, returns one box per object & JSON & NumPy, SciPy & voxel size, distance & 6--40 & 0.2--16 & 199 \\
        \textbf{normals} & estimates the surface normal of every point & normals & NumPy, SciPy & neighbors & 6--40 & 1.8--46 & 198 \\
        \textbf{register} & aligns a scan with the previous one (rotation, translation) & JSON & NumPy, SciPy & iterations & 6--38 & 1.0--45 & 189 \\
        \textbf{birdview} & renders a bird's-eye-view image of the scan & image & NumPy & -- & 6--40 & 1.1--3.2 & 197 \\
        \textbf{encode} & quantizes and packs the point coordinates & bytes & NumPy & precision & 6--38 & 3.2--18 & 200 \\
        \midrule
        \multicolumn{8}{@{}l}{\textit{Vision} (data quality: frame height, pixels; for object and face also the model size)} \\
        \textbf{object} & detects objects (YOLOv10): boxes, classes, scores & JSON & ONNX Runtime & -- & 160--640; 1--2 & 13--320 & 182 \\
        \textbf{face} & detects faces (YuNet): boxes and scores & JSON & OpenCV & -- & 256--640; 1--2 & 7.6--195 & 171 \\
        \textbf{edge} & finds edges (bilateral filter, Canny) and their contours & JSON & OpenCV & smoothing & 160--640 & 0.9--6.2 & 200 \\
        \textbf{motion} & returns the regions moved since the previous frame & JSON & OpenCV & method & 160--640 & 0.8--6.0 & 200 \\
        \textbf{QR} & finds and decodes QR codes, draws them onto the frame & image & OpenCV, pyzbar & -- & 192--608 & 1.5--9.1 & 200 \\
        \textbf{transform} & resizes a frame, corrects its colors and geometry & image & OpenCV & correction & 192--640 & 7.1--41 & 200 \\
        \textbf{compress} & encodes a frame as a compressed image & bytes & OpenCV, Pillow & quality, format & 160--608 & 0.2--1.0 & 181 \\
        \textbf{denoise} & removes noise, keeps edges (bilateral filter) & image & OpenCV & strength & 160--640 & 1.4--22 & 198 \\
        \bottomrule
    \end{tabular}
\end{center}
\vspace{6pt}
}]
\makeatother

\noindent
\textbf{Services.}
%
Items are LibriSpeech \textit{dev-clean} clips of at least 12\,s\footnote{\url{https://www.openslr.org/12}}, scans of the KITTI raw recording \texttt{2011\_09\_26\_drive\_0001\_sync}\footnote{\url{https://www.cvlibs.net/datasets/kitti/raw_data.php}}, and frames of sample videos. The data quality knob sets the length of the processed clip, the radius around the car within which points are kept, and the height to which a frame is resized. Three services run pretrained models: transcribe (Vosk\footnote{\url{https://alphacephei.com/vosk/models}}), object (YOLOv10\footnote{\url{https://github.com/THU-MIG/yolov10}}), and face (YuNet\footnote{\url{https://github.com/opencv/opencv_zoo/tree/main/models/face_detection_yunet}}). For object and face, the model size is a second data quality knob: 1 selects a small model, 2 a larger one, which takes several times as long per frame. All other knobs stay at their default.

\vspace{3pt}
\noindent
\textbf{Benchmarks.}
%
Every service runs alone with $c = 1, \dots, 4$ workers in the same 10 configurations, each at the utilizations 0.1, 0.4, 0.6, 0.8, and 1.0 per worker. A run has the CPU need of its configuration, measured in an 8\,s probe, plus 15\% as cores, and measures 30\,s or 400 items (at most 180\,s) after a 10\,s warm-up. Invalid runs, in which the service fell behind its arrivals or needed more cores than its probe had, are dropped (\textit{Runs}). The GPs of a service take as inputs its data quality (and model size), its workers $c$, and the logarithm of its arrival rate in items/s; the cores are no input, since no used run is throttled. Where a benchmark covers less data quality than the compositions (e.g., face from 256\,px), the models extrapolate.

\vspace{3pt}
\noindent
\textbf{Compositions.}
%
Each composition is measured in 200 runs with the measurement window of the benchmarks: 150 runs with one worker and 1, 1.5, or 2 cores per service, and 50 runs with $c \in \{1, \dots, 4\}$ workers and $c + 0.5$ cores per service. The data quality is drawn in steps of 1\,s (Audio), 2\,m (Lidar), and 32\,px (Vision). Only runs in which no service was throttled are used: 125--156 per composition (75--106 with one worker, all 50 with several).

\vspace{3pt}
\noindent
\textbf{Example.}
%
% Source: statics/results_milos_v2/b1/b1-vision-5-single/*/runs/run-0000 (trace_steps.csv, <node>/records.csv, run.json);
% the item motion-376, the median end-to-end latency among the 415 items of the measurement phase.
Table~\ref{tab:apx-trace} follows one item through the six services of V5 in a run with 576\,px, model size 1, and one worker per service; its end-to-end latency is the median of the 415 items of that run. Four services return JSON, QR the annotated frame as a raw image, and compress a JPEG (quality 75). The 157\,ms of the item are the sum of the times in the table; processing at the object detector alone takes two thirds of it.

\makeatletter
\begin{center}
    \def\@captype{table}
    \caption{One item of V5 per step (576\,px, model size 1): output and its size, time waited in the queue, handed over from the previous service, and processed, in ms.}
    \label{tab:apx-trace}
    \footnotesize
    \setlength{\tabcolsep}{4pt}
    \begin{tabular}{@{}l l r r r r@{}}
        \toprule
        {Service} & {Output} & {Size} & {Wait} & {Handover} & {Proc.} \\
        \midrule
        \textbf{motion} & JSON (moved regions) & -- & 0.3 & -- & 5.6 \\
        \textbf{QR} & raw image (annotated) & 1.77\,MB & 0.4 & 1.9 & 8.1 \\
        \textbf{object} & JSON (detections) & -- & 3.8 & 1.3 & 102.7 \\
        \textbf{face} & JSON (faces) & -- & 0.4 & 1.5 & 21.2 \\
        \textbf{edge} & JSON (contours) & -- & 0.4 & 1.5 & 5.6 \\
        \textbf{compress} & JPEG & 38.8\,kB & 0.4 & 1.3 & 1.0 \\
        \midrule
        \multicolumn{5}{@{}l}{end-to-end latency} & 157.3 \\
        \bottomrule
    \end{tabular}
\end{center}
\makeatother

% \twocolumn[...] starts a new page and sets the heading and the service table across both columns at its top;
% the text below continues in two columns. The table is no float (a table* would be pushed to the page after),
% so \@captype gives it a caption and number of a regular table.
% Service names and the order of the applications (Audio, Lidar, Vision) follow Table~\ref{tab:applications}.
% The text only adds what Sec. V does not state; it must fit on one page together with the table.
% The footnotes use \url (packages url and hyperref in main.tex).

\makeatletter
\twocolumn[{%
\section{Detailed Description of Services and Compositions}
\label{apx:services}
\begin{center}
    \def\@captype{table}
    \caption{Services per application. \textit{Fixed knobs} stay at their default in all runs; \textit{Data Q.}\ is the data quality range covered by the benchmarks, \textit{Proc.}\ the range of the median processing time of an item over the benchmarked configurations with one worker, and \textit{Runs} the number of valid benchmark runs with 1--4 workers.}
    \label{tab:apx-services}
    \scriptsize
    \setlength{\tabcolsep}{4pt}
    \begin{tabular}{@{}l p{0.33\textwidth} l l l l r r@{}}
        \toprule
        {Service} & {Function} & {Output} & {Built on} & {Fixed knobs} & {Data Q.} & {Proc.\ [ms]} & {Runs} \\
        \midrule
        \multicolumn{8}{@{}l}{\textit{Audio} (data quality: clip length, seconds)} \\
        \textbf{resample} & resamples a clip, removes its offset, normalizes loudness & audio & NumPy, SciPy & target rate & 2--12 & 0.6--3.9 & 200 \\
        \textbf{denoise} & suppresses background noise in the spectrum & audio & NumPy & FFT size & 2--12 & 6.4--35 & 194 \\
        \textbf{MFCC} & computes a mel spectrogram and its MFCC features & features & NumPy & mel bands, frame & 2--12 & 0.9--4.7 & 200 \\
        \textbf{voicedetect} & returns the start and end times of speech segments & JSON & NumPy & frame length & 3--12 & 4.3--16 & 200 \\
        \textbf{pitch} & estimates the pitch of the voice per frame & pitch & NumPy & hop length & 2--12 & 7.6--31 & 187 \\
        \textbf{transcribe} & transcribes speech to text (small model) & text & Vosk & model size & 2--12 & 337--983 & 164 \\
        \textbf{encode} & encodes a clip as a compressed audio file & bytes & NumPy, soundfile & format & 2--12 & 1.6--9.1 & 200 \\
        \midrule
        \multicolumn{8}{@{}l}{\textit{Lidar} (data quality: max.\ distance, meters)} \\
        \textbf{voxel} & downsamples a scan to one point per voxel & points & NumPy & voxel size & 8--40 & 1.9--9.0 & 197 \\
        \textbf{ground} & fits the ground plane by random trials, removes it & points & NumPy & attempts & 8--40 & 17--47 & 194 \\
        \textbf{cluster} & clusters the points, returns one box per object & JSON & NumPy, SciPy & voxel size, distance & 6--40 & 0.2--16 & 199 \\
        \textbf{normals} & estimates the surface normal of every point & normals & NumPy, SciPy & neighbors & 6--40 & 1.8--46 & 198 \\
        \textbf{register} & aligns a scan with the previous one (rotation, translation) & JSON & NumPy, SciPy & iterations & 6--38 & 1.0--45 & 189 \\
        \textbf{birdview} & renders a bird's-eye-view image of the scan & image & NumPy & -- & 6--40 & 1.1--3.2 & 197 \\
        \textbf{encode} & quantizes and packs the point coordinates & bytes & NumPy & precision & 6--38 & 3.2--18 & 200 \\
        \midrule
        \multicolumn{8}{@{}l}{\textit{Vision} (data quality: frame height, pixels; for object and face also the model size)} \\
        \textbf{object} & detects objects (YOLOv10): boxes, classes, scores & JSON & ONNX Runtime & -- & 160--640; 1--2 & 13--320 & 182 \\
        \textbf{face} & detects faces (YuNet): boxes and scores & JSON & OpenCV & -- & 256--640; 1--2 & 7.6--195 & 171 \\
        \textbf{edge} & finds edges (bilateral filter, Canny) and their contours & JSON & OpenCV & smoothing & 160--640 & 0.9--6.2 & 200 \\
        \textbf{motion} & returns the regions moved since the previous frame & JSON & OpenCV & method & 160--640 & 0.8--6.0 & 200 \\
        \textbf{QR} & finds and decodes QR codes, draws them onto the frame & image & OpenCV, pyzbar & -- & 192--608 & 1.5--9.1 & 200 \\
        \textbf{transform} & resizes a frame, corrects its colors and geometry & image & OpenCV & correction & 192--640 & 7.1--41 & 200 \\
        \textbf{compress} & encodes a frame as a compressed image & bytes & OpenCV, Pillow & quality, format & 160--608 & 0.2--1.0 & 181 \\
        \textbf{denoise} & removes noise, keeps edges (bilateral filter) & image & OpenCV & strength & 160--640 & 1.4--22 & 198 \\
        \bottomrule
    \end{tabular}
\end{center}
\vspace{6pt}
}]
\makeatother

\noindent
\textbf{Services.}
%
Items are LibriSpeech \textit{dev-clean} clips of at least 12\,s\footnote{\url{https://www.openslr.org/12}}, scans of the KITTI raw recording \texttt{2011\_09\_26\_drive\_0001\_sync}\footnote{\url{https://www.cvlibs.net/datasets/kitti/raw_data.php}}, and frames of sample videos. The data quality knob sets the length of the processed clip, the radius around the car within which points are kept, and the height to which a frame is resized. Three services run pretrained models: transcribe (Vosk\footnote{\url{https://alphacephei.com/vosk/models}}), object (YOLOv10\footnote{\url{https://github.com/THU-MIG/yolov10}}), and face (YuNet\footnote{\url{https://github.com/opencv/opencv_zoo/tree/main/models/face_detection_yunet}}). For object and face, the model size is a second data quality knob: 1 selects a small model, 2 a larger one, which takes several times as long per frame. All other knobs stay at their default.

\vspace{3pt}
\noindent
\textbf{Benchmarks.}
%
Every service runs alone with $c = 1, \dots, 4$ workers in the same 10 configurations, each at the utilizations 0.1, 0.4, 0.6, 0.8, and 1.0 per worker. A run has the CPU need of its configuration, measured in an 8\,s probe, plus 15\% as cores, and measures 30\,s or 400 items (at most 180\,s) after a 10\,s warm-up. Invalid runs, in which the service fell behind its arrivals or needed more cores than its probe had, are dropped (\textit{Runs}). The GPs of a service take as inputs its data quality (and model size), its workers $c$, and the logarithm of its arrival rate in items/s; the cores are no input, since no used run is throttled. Where a benchmark covers less data quality than the compositions (e.g., face from 256\,px), the models extrapolate.

\vspace{3pt}
\noindent
\textbf{Compositions.}
%
Each composition is measured in 200 runs with the measurement window of the benchmarks: 150 runs with one worker and 1, 1.5, or 2 cores per service, and 50 runs with $c \in \{1, \dots, 4\}$ workers and $c + 0.5$ cores per service. The data quality is drawn in steps of 1\,s (Audio), 2\,m (Lidar), and 32\,px (Vision). Only runs in which no service was throttled are used: 125--156 per composition (75--106 with one worker, all 50 with several).

\vspace{3pt}
\noindent
\textbf{Example.}
%
% Source: statics/results_milos_v2/b1/b1-vision-5-single/*/runs/run-0000 (trace_steps.csv, <node>/records.csv, run.json);
% the item motion-376, the median end-to-end latency among the 415 items of the measurement phase.
Table~\ref{tab:apx-trace} follows one item through the six services of V5 in a run with 576\,px, model size 1, and one worker per service; its end-to-end latency is the median of the 415 items of that run. Four services return JSON, QR the annotated frame as a raw image, and compress a JPEG (quality 75). The 157\,ms of the item are the sum of the times in the table; processing at the object detector alone takes two thirds of it.

\makeatletter
\begin{center}
    \def\@captype{table}
    \caption{One item of V5 per step (576\,px, model size 1): output and its size, time waited in the queue, handed over from the previous service, and processed, in ms.}
    \label{tab:apx-trace}
    \footnotesize
    \setlength{\tabcolsep}{4pt}
    \begin{tabular}{@{}l l r r r r@{}}
        \toprule
        {Service} & {Output} & {Size} & {Wait} & {Handover} & {Proc.} \\
        \midrule
        \textbf{motion} & JSON (moved regions) & -- & 0.3 & -- & 5.6 \\
        \textbf{QR} & raw image (annotated) & 1.77\,MB & 0.4 & 1.9 & 8.1 \\
        \textbf{object} & JSON (detections) & -- & 3.8 & 1.3 & 102.7 \\
        \textbf{face} & JSON (faces) & -- & 0.4 & 1.5 & 21.2 \\
        \textbf{edge} & JSON (contours) & -- & 0.4 & 1.5 & 5.6 \\
        \textbf{compress} & JPEG & 38.8\,kB & 0.4 & 1.3 & 1.0 \\
        \midrule
        \multicolumn{5}{@{}l}{end-to-end latency} & 157.3 \\
        \bottomrule
    \end{tabular}
\end{center}
\makeatother

% \twocolumn[...] starts a new page and sets the heading and the service table across both columns at its top;
% the text below continues in two columns. The table is no float (a table* would be pushed to the page after),
% so \@captype gives it a caption and number of a regular table.
% Service names and the order of the applications (Audio, Lidar, Vision) follow Table~\ref{tab:applications}.
% The text only adds what Sec. V does not state; it must fit on one page together with the table.
% The footnotes use \url (packages url and hyperref in main.tex).

\makeatletter
\twocolumn[{%
\section{Detailed Description of Services and Compositions}
\label{apx:services}
\begin{center}
    \def\@captype{table}
    \caption{Services per application. \textit{Fixed knobs} stay at their default in all runs; \textit{Data Q.}\ is the data quality range covered by the benchmarks, \textit{Proc.}\ the range of the median processing time of an item over the benchmarked configurations with one worker, and \textit{Runs} the number of valid benchmark runs with 1--4 workers.}
    \label{tab:apx-services}
    \scriptsize
    \setlength{\tabcolsep}{4pt}
    \begin{tabular}{@{}l p{0.33\textwidth} l l l l r r@{}}
        \toprule
        {Service} & {Function} & {Output} & {Built on} & {Fixed knobs} & {Data Q.} & {Proc.\ [ms]} & {Runs} \\
        \midrule
        \multicolumn{8}{@{}l}{\textit{Audio} (data quality: clip length, seconds)} \\
        \textbf{resample} & resamples a clip, removes its offset, normalizes loudness & audio & NumPy, SciPy & target rate & 2--12 & 0.6--3.9 & 200 \\
        \textbf{denoise} & suppresses background noise in the spectrum & audio & NumPy & FFT size & 2--12 & 6.4--35 & 194 \\
        \textbf{MFCC} & computes a mel spectrogram and its MFCC features & features & NumPy & mel bands, frame & 2--12 & 0.9--4.7 & 200 \\
        \textbf{voicedetect} & returns the start and end times of speech segments & JSON & NumPy & frame length & 3--12 & 4.3--16 & 200 \\
        \textbf{pitch} & estimates the pitch of the voice per frame & pitch & NumPy & hop length & 2--12 & 7.6--31 & 187 \\
        \textbf{transcribe} & transcribes speech to text (small model) & text & Vosk & model size & 2--12 & 337--983 & 164 \\
        \textbf{encode} & encodes a clip as a compressed audio file & bytes & NumPy, soundfile & format & 2--12 & 1.6--9.1 & 200 \\
        \midrule
        \multicolumn{8}{@{}l}{\textit{Lidar} (data quality: max.\ distance, meters)} \\
        \textbf{voxel} & downsamples a scan to one point per voxel & points & NumPy & voxel size & 8--40 & 1.9--9.0 & 197 \\
        \textbf{ground} & fits the ground plane by random trials, removes it & points & NumPy & attempts & 8--40 & 17--47 & 194 \\
        \textbf{cluster} & clusters the points, returns one box per object & JSON & NumPy, SciPy & voxel size, distance & 6--40 & 0.2--16 & 199 \\
        \textbf{normals} & estimates the surface normal of every point & normals & NumPy, SciPy & neighbors & 6--40 & 1.8--46 & 198 \\
        \textbf{register} & aligns a scan with the previous one (rotation, translation) & JSON & NumPy, SciPy & iterations & 6--38 & 1.0--45 & 189 \\
        \textbf{birdview} & renders a bird's-eye-view image of the scan & image & NumPy & -- & 6--40 & 1.1--3.2 & 197 \\
        \textbf{encode} & quantizes and packs the point coordinates & bytes & NumPy & precision & 6--38 & 3.2--18 & 200 \\
        \midrule
        \multicolumn{8}{@{}l}{\textit{Vision} (data quality: frame height, pixels; for object and face also the model size)} \\
        \textbf{object} & detects objects (YOLOv10): boxes, classes, scores & JSON & ONNX Runtime & -- & 160--640; 1--2 & 13--320 & 182 \\
        \textbf{face} & detects faces (YuNet): boxes and scores & JSON & OpenCV & -- & 256--640; 1--2 & 7.6--195 & 171 \\
        \textbf{edge} & finds edges (bilateral filter, Canny) and their contours & JSON & OpenCV & smoothing & 160--640 & 0.9--6.2 & 200 \\
        \textbf{motion} & returns the regions moved since the previous frame & JSON & OpenCV & method & 160--640 & 0.8--6.0 & 200 \\
        \textbf{QR} & finds and decodes QR codes, draws them onto the frame & image & OpenCV, pyzbar & -- & 192--608 & 1.5--9.1 & 200 \\
        \textbf{transform} & resizes a frame, corrects its colors and geometry & image & OpenCV & correction & 192--640 & 7.1--41 & 200 \\
        \textbf{compress} & encodes a frame as a compressed image & bytes & OpenCV, Pillow & quality, format & 160--608 & 0.2--1.0 & 181 \\
        \textbf{denoise} & removes noise, keeps edges (bilateral filter) & image & OpenCV & strength & 160--640 & 1.4--22 & 198 \\
        \bottomrule
    \end{tabular}
\end{center}
\vspace{6pt}
}]
\makeatother

\noindent
\textbf{Services.}
%
Items are LibriSpeech \textit{dev-clean} clips of at least 12\,s\footnote{\url{https://www.openslr.org/12}}, scans of the KITTI raw recording \texttt{2011\_09\_26\_drive\_0001\_sync}\footnote{\url{https://www.cvlibs.net/datasets/kitti/raw_data.php}}, and frames of sample videos. The data quality knob sets the length of the processed clip, the radius around the car within which points are kept, and the height to which a frame is resized. Three services run pretrained models: transcribe (Vosk\footnote{\url{https://alphacephei.com/vosk/models}}), object (YOLOv10\footnote{\url{https://github.com/THU-MIG/yolov10}}), and face (YuNet\footnote{\url{https://github.com/opencv/opencv_zoo/tree/main/models/face_detection_yunet}}). For object and face, the model size is a second data quality knob: 1 selects a small model, 2 a larger one, which takes several times as long per frame. All other knobs stay at their default.

\vspace{3pt}
\noindent
\textbf{Benchmarks.}
%
Every service runs alone with $c = 1, \dots, 4$ workers in the same 10 configurations, each at the utilizations 0.1, 0.4, 0.6, 0.8, and 1.0 per worker. A run has the CPU need of its configuration, measured in an 8\,s probe, plus 15\% as cores, and measures 30\,s or 400 items (at most 180\,s) after a 10\,s warm-up. Invalid runs, in which the service fell behind its arrivals or needed more cores than its probe had, are dropped (\textit{Runs}). The GPs of a service take as inputs its data quality (and model size), its workers $c$, and the logarithm of its arrival rate in items/s; the cores are no input, since no used run is throttled. Where a benchmark covers less data quality than the compositions (e.g., face from 256\,px), the models extrapolate.

\vspace{3pt}
\noindent
\textbf{Compositions.}
%
Each composition is measured in 200 runs with the measurement window of the benchmarks: 150 runs with one worker and 1, 1.5, or 2 cores per service, and 50 runs with $c \in \{1, \dots, 4\}$ workers and $c + 0.5$ cores per service. The data quality is drawn in steps of 1\,s (Audio), 2\,m (Lidar), and 32\,px (Vision). Only runs in which no service was throttled are used: 125--156 per composition (75--106 with one worker, all 50 with several).

\vspace{3pt}
\noindent
\textbf{Example.}
%
% Source: statics/results_milos_v2/b1/b1-vision-5-single/*/runs/run-0000 (trace_steps.csv, <node>/records.csv, run.json);
% the item motion-376, the median end-to-end latency among the 415 items of the measurement phase.
Table~\ref{tab:apx-trace} follows one item through the six services of V5 in a run with 576\,px, model size 1, and one worker per service; its end-to-end latency is the median of the 415 items of that run. Four services return JSON, QR the annotated frame as a raw image, and compress a JPEG (quality 75). The 157\,ms of the item are the sum of the times in the table; processing at the object detector alone takes two thirds of it.

\makeatletter
\begin{center}
    \def\@captype{table}
    \caption{One item of V5 per step (576\,px, model size 1): output and its size, time waited in the queue, handed over from the previous service, and processed, in ms.}
    \label{tab:apx-trace}
    \footnotesize
    \setlength{\tabcolsep}{4pt}
    \begin{tabular}{@{}l l r r r r@{}}
        \toprule
        {Service} & {Output} & {Size} & {Wait} & {Handover} & {Proc.} \\
        \midrule
        \textbf{motion} & JSON (moved regions) & -- & 0.3 & -- & 5.6 \\
        \textbf{QR} & raw image (annotated) & 1.77\,MB & 0.4 & 1.9 & 8.1 \\
        \textbf{object} & JSON (detections) & -- & 3.8 & 1.3 & 102.7 \\
        \textbf{face} & JSON (faces) & -- & 0.4 & 1.5 & 21.2 \\
        \textbf{edge} & JSON (contours) & -- & 0.4 & 1.5 & 5.6 \\
        \textbf{compress} & JPEG & 38.8\,kB & 0.4 & 1.3 & 1.0 \\
        \midrule
        \multicolumn{5}{@{}l}{end-to-end latency} & 157.3 \\
        \bottomrule
    \end{tabular}
\end{center}
\makeatother
